\documentclass{situationist-preprint}

\title{Bohmian Mechanics as an Empirically Equivalent Special Case}
\msid{SCI-09-c}
\msfamily{SCI-09}
\msdiscipline{Philosophy of Science}
\mstype{Alternative proof}
\msstatus{Revised 7 October 2026}
\mskeywords{Bohmian mechanics; Everett interpretation; empirical equivalence; quantum equilibrium; underdetermination}

\begin{document}
\maketitle

\begin{abstract}
Bohmian mechanics adds to the wave function a configuration of particles guided by it. In quantum equilibrium, it is empirically equivalent to standard quantum mechanics and therefore to the Everett interpretation defended in SCI-09-a. We argue that this equivalence supports Everett: Bohmian mechanics contains the full Everettian wave function as a dynamical object, and the particle configuration adds nothing that affects predictions in equilibrium.
\end{abstract}

\paragraph{Note on the revision.} The version of 6 October 2026 described Bohmian mechanics as an approximation to the Everett interpretation. This is incorrect: in quantum equilibrium the two are empirically equivalent, not related by approximation. The title and Section~1 have been corrected.

\section{Equivalence}

In Bohmian mechanics, the state of a system consists of a wave function, evolving according to the Schrödinger equation, and an actual configuration of particles, moving according to a guidance equation determined by the wave function \citep{Bohm1952}. If the particles are initially distributed according to the Born rule, a condition called quantum equilibrium, they remain so distributed, and the theory reproduces the statistical predictions of standard quantum mechanics \citep{DurrGoldsteinZanghi1992}.

\begin{proposition}\label{p:eq}
In quantum equilibrium, Bohmian mechanics and the Everett interpretation make the same empirical predictions.
\end{proposition}

\begin{proof}
Both theories evolve the wave function by the Schrödinger equation without collapse. In Bohmian mechanics, observed frequencies are given by the equilibrium distribution of configurations, which is the Born distribution. In the Everett interpretation, the Born rule is derived from decision-theoretic considerations as the rational weighting of branches (SCI-09-a). The two theories therefore assign the same probabilities to measurement outcomes.
\end{proof}

Outside equilibrium the theories could in principle diverge \citep{Valentini1991}, but no evidence of non-equilibrium has been found, and we confine the comparison to the equilibrium regime.

\section{Why equivalence favours Everett}

Empirical equivalence does not leave the choice open. The Bohmian wave function is the same object as the Everettian one: it evolves in the same way and has the same branching structure. The particle configuration selects one branch as the occupied one, but the other branches remain in the wave function and have the same structure as the occupied one. As \citet{Deutsch1996} put it, Bohmian mechanics is a theory of parallel universes in a state of chronic denial; \citet{BrownWallace2005} develop the point in detail. Given empirical equivalence, the Everett interpretation is the simpler theory, since it postulates only the wave function, and Bohmian mechanics is the special case in which an additional variable is attached to one branch.

\begin{thebibliography}{9}\small
\bibitem[Bohm(1952)]{Bohm1952} Bohm, D. (1952). A suggested interpretation of the quantum theory in terms of ``hidden'' variables, I and II. \emph{Physical Review}, 85(2), 166--193.
\bibitem[Brown and Wallace(2005)]{BrownWallace2005} Brown, H.~R. and Wallace, D. (2005). Solving the measurement problem: De Broglie--Bohm loses out to Everett. \emph{Foundations of Physics}, 35(4), 517--540.
\bibitem[Deutsch(1996)]{Deutsch1996} Deutsch, D. (1996). Comment on Lockwood. \emph{The British Journal for the Philosophy of Science}, 47(2), 222--228.
\bibitem[Dürr et~al.(1992)]{DurrGoldsteinZanghi1992} Dürr, D., Goldstein, S. and Zanghì, N. (1992). Quantum equilibrium and the origin of absolute uncertainty. \emph{Journal of Statistical Physics}, 67(5--6), 843--907.
\bibitem[Valentini(1991)]{Valentini1991} Valentini, A. (1991). Signal-locality, uncertainty, and the subquantum H-theorem, I. \emph{Physics Letters A}, 156(1--2), 5--11.
\end{thebibliography}

\end{document}
